\documentclass{article}
\usepackage{iclr2026_conference,times}

\usepackage{amsmath,amsfonts,amssymb}
\usepackage{graphicx}
\usepackage{booktabs}
\usepackage{algorithm}
<parameter name="content">\usepackage{algorithmic}

\title{Hierarchical Reinforcement Learning for Robotic Manipulation with Nullspace Optimization}

\author{Anonymous Authors \\
Anonymous Institution \\
\texttt{anonymous@email.com}}

\begin{document}

\maketitle

\begin{abstract}
机械臂操作任务需要同时满足任务目标和物理约束。本文提出一种结合零空间优化的分层强化学习方法，在任务空间学习高层策略的同时，利用零空间投影优化关节空间的物理性能。实验表明该方法在复杂场景下显著提升成功率和样本效率。
\end{abstract}

\section{Introduction}

机械臂操作是机器人领域的核心问题。传统方法依赖精确建模，难以适应复杂环境；纯强化学习方法样本效率低且难以保证物理可行性。

本文贡献：
\begin{itemize}
\item 提出分层控制架构，解耦任务目标与物理约束
\item 设计零空间优化模块，提升关节空间性能
\item 在多种场景下验证方法有效性
\end{itemize}

\section{Related Work}

\subsection{传统运动规划}
RRT*等方法需要精确环境模型，计算复杂度高。

\subsection{强化学习方法}
PPO、SAC等方法样本效率低，难以处理高维连续控制。

\subsection{分层控制}
现有方法缺乏对冗余自由度的显式优化。

\section{Method}

\subsection{问题定义}

\textbf{任务目标：}
\begin{itemize}
\item 末端到达目标位置
\item 避障
\item 关节力矩平滑
\end{itemize}

\textbf{控制结构：}采用分层架构，高层策略输出任务空间速度$\dot{x}_d \in \mathbb{R}^6$，低层控制器映射到关节空间。

\textbf{强化学习形式化：}MDP定义为$(\mathcal{S}, \mathcal{A}, \mathcal{P}, \mathcal{R}, \gamma)$，其中：
\begin{itemize}
\item 状态$s_t$：末端位姿、目标位置、障碍物信息
\item 动作$a_t$：任务空间速度增量
\item 奖励$r_t$：到达奖励 - 碰撞惩罚 - 力矩惩罚
\end{itemize}

\subsection{机械臂动力学模型}

\textbf{运动学关系：}
\begin{equation}
\dot{x} = J(q)\dot{q}
\end{equation}
其中$J(q) \in \mathbb{R}^{6 \times 7}$为雅可比矩阵。

\textbf{动力学方程：}
\begin{equation}
M(q)\ddot{q} + C(q,\dot{q})\dot{q} + g(q) = \tau
\end{equation}

\textbf{零空间投影：}
\begin{equation}
N(q) = I - J^\dagger(q)J(q)
\end{equation}

\subsection{零空间优化模块}

\textbf{动机：}7自由度机械臂存在冗余，可在完成任务的同时优化次要目标。

\textbf{理论依据：}零空间投影保证不影响末端运动，同时优化关节配置。

\textbf{数学形式化：}
\begin{equation}
\dot{q} = J^\dagger(q)\dot{x}_d + N(q)\dot{q}_0
\end{equation}
其中$\dot{q}_0$为次要任务速度（如远离关节限位）。

\textbf{与baseline对比：}相比纯reward shaping，零空间优化保持最优性且训练更稳定。

\section{Experiments}

\subsection{实验设置}

\textbf{对比方法：}
\begin{itemize}
\item RRT*
\item CHOMP
\item PPO-Baseline
\item SAC-Baseline
\item Ours-NoNullspace（消融）
\end{itemize}

\textbf{测试场景：}
\begin{itemize}
\item 场景1：稀疏障碍（3个静态障碍物）
\item 场景2：密集障碍（12个障碍物）
\item 场景3：动态障碍（2个运动障碍物）
\end{itemize}

\textbf{评价指标：}
\begin{itemize}
\item 任务成功率
\item 样本效率（训练步数）
\item 路径长度
\item 关节力矩平滑度
\item 最小障碍距离
\end{itemize}

\subsection{结果分析}

实验结果表明，本文方法在所有场景下均优于baseline，特别是在密集障碍场景下成功率提升显著。

\section{Conclusion}

本文提出的分层强化学习方法有效结合了任务空间学习与零空间优化，在复杂操作任务中展现出优越性能。

\bibliographystyle{iclr2026_conference}
\bibliography{references}

\end{document}
